\begin{proof}[Proof of \cref{obj:lem:randomized-kernel-tv-comparison}]
\begin{enumerate}
\item We first establish the bounded-function estimate used for the kernel probabilities. For any measurable function
\[
\varphi:\mathcal Z\longrightarrow[0,1],
\qquad
L_{\varphi,t}:=\{z\in\mathcal Z:t\le\varphi(z)\}
\quad(t\in\mathbb R),
\]
the level sets are measurable, and the definition of total variation gives
\[
\bigl|Q_0(L_{\varphi,t})-Q_1(L_{\varphi,t})\bigr|
\le \operatorname{TV}(Q_0,Q_1).
\]
The layer-cake identity gives, for \(i\in\{0,1\}\),
\[
\int_{\mathcal Z}\varphi(z)\,Q_i(dz)
=
\int_{(0,1]}Q_i(L_{\varphi,t})\,dt.
\]
All integrands are measurable and bounded. Subtracting these identities and integrating the preceding event bound yields
\[
\begin{aligned}
\left|
\int_{\mathcal Z}\varphi\,dQ_0-
\int_{\mathcal Z}\varphi\,dQ_1
\right|
&=
\left|\int_{(0,1]}
\bigl[Q_0(L_{\varphi,t})-Q_1(L_{\varphi,t})\bigr]\,dt\right|\\
&\le
\int_{(0,1]}
\bigl|Q_0(L_{\varphi,t})-Q_1(L_{\varphi,t})\bigr|\,dt\\
&\le \operatorname{TV}(Q_0,Q_1).
\end{aligned}
\]
% lean: Causalean.Stat.tvDist_integral_range
For any measurable event \(T_{\mathrm{act}}\subseteq\mathcal H\), define
\[
\varphi_{\mathrm{act}}(z):=\kappa(z,T_{\mathrm{act}})
\quad(z\in\mathcal Z).
\]
This function is measurable and takes values in \([0,1]\), and
\[
(Q_i\kappa)(T_{\mathrm{act}})
=
\int_{\mathcal Z}\varphi_{\mathrm{act}}\,dQ_i
\quad(i\in\{0,1\}).
\]
Applying the bounded-function estimate and taking the supremum over measurable action events proves
\[
\operatorname{TV}(Q_0\kappa,Q_1\kappa)
\le \operatorname{TV}(Q_0,Q_1).
\]
% lean: Causalean.Stat.tvDist_bind_le

\item For any measurable joint event \(T_{\mathrm{joint}}\subseteq\mathcal Z\times\mathcal H\), define its sections and their kernel probabilities by
\[
T_{\mathrm{joint},z}
:=
\{h\in\mathcal H:(z,h)\in T_{\mathrm{joint}}\},
\qquad
\varphi_{\mathrm{joint}}(z)
:=
\kappa(z,T_{\mathrm{joint},z})
\quad(z\in\mathcal Z).
\]
The sections are measurable, and \(\varphi_{\mathrm{joint}}\) is measurable with values in \([0,1]\). The joint-law identity is
\[
(Q_i\ltimes\kappa)(T_{\mathrm{joint}})
=
\int_{\mathcal Z}\varphi_{\mathrm{joint}}\,dQ_i
\quad(i\in\{0,1\}).
\]
Thus the bounded-function estimate, followed by the supremum over joint events, gives
\[
\operatorname{TV}(Q_0\ltimes\kappa,Q_1\ltimes\kappa)
\le \operatorname{TV}(Q_0,Q_1).
\]
For the reverse inequality, every measurable input event \(T_{\mathrm{in}}\subseteq\mathcal Z\) satisfies
\[
(Q_i\ltimes\kappa)(T_{\mathrm{in}}\times\mathcal H)
=
Q_i(T_{\mathrm{in}})
\quad(i\in\{0,1\}),
\]
because \(\kappa(z,\mathcal H)=1\). Consequently,
\[
\begin{aligned}
|Q_0(T_{\mathrm{in}})-Q_1(T_{\mathrm{in}})|
&=
\left|
(Q_0\ltimes\kappa)(T_{\mathrm{in}}\times\mathcal H)
-
(Q_1\ltimes\kappa)(T_{\mathrm{in}}\times\mathcal H)
\right|\\
&\le
\operatorname{TV}(Q_0\ltimes\kappa,Q_1\ltimes\kappa).
\end{aligned}
\]
Taking the supremum over input events proves
\[
\operatorname{TV}(Q_0\ltimes\kappa,Q_1\ltimes\kappa)
=
\operatorname{TV}(Q_0,Q_1).
\]
% lean: Causalean.Stat.tvDist_compProd_eq
Combining this equality with the action-marginal bound gives
\[
\operatorname{TV}(Q_0\kappa,Q_1\kappa)
\le
\operatorname{TV}(Q_0\ltimes\kappa,Q_1\ltimes\kappa)
=
\operatorname{TV}(Q_0,Q_1).
\]
% lean: Causalean.Stat.tvDist_bind_le_compProd

\item The joint laws are probability laws. Hence, for any measurable joint event \(T_{\mathrm{joint}}\),
\[
(Q_1\ltimes\kappa)(T_{\mathrm{joint}}^c)
=
1-(Q_1\ltimes\kappa)(T_{\mathrm{joint}}),
\]
and the event bound defining their total variation gives
\[
\begin{aligned}
&(Q_0\ltimes\kappa)(T_{\mathrm{joint}})
 +(Q_1\ltimes\kappa)(T_{\mathrm{joint}}^c)\\
&\qquad=
1+(Q_0\ltimes\kappa)(T_{\mathrm{joint}})
 -(Q_1\ltimes\kappa)(T_{\mathrm{joint}})\\
&\qquad\ge
1-\operatorname{TV}(Q_0\ltimes\kappa,Q_1\ltimes\kappa)
=
1-\operatorname{TV}(Q_0,Q_1).
\end{aligned}
\]
% lean: Causalean.Stat.one_sub_tvDist_le_compProd_test

\item Likewise, the action marginals are probability laws, so every measurable action event \(T_{\mathrm{act}}\) satisfies
\[
(Q_1\kappa)(T_{\mathrm{act}}^c)
=
1-(Q_1\kappa)(T_{\mathrm{act}}).
\]
Using their total-variation event bound and then the action-marginal contraction established above,
\[
\begin{aligned}
(Q_0\kappa)(T_{\mathrm{act}})
 +(Q_1\kappa)(T_{\mathrm{act}}^c)
&=
1+(Q_0\kappa)(T_{\mathrm{act}})
 -(Q_1\kappa)(T_{\mathrm{act}})\\
&\ge
1-\operatorname{TV}(Q_0\kappa,Q_1\kappa)\\
&\ge
1-\operatorname{TV}(Q_0,Q_1).
\end{aligned}
\]
% lean: Causalean.Stat.one_sub_tvDist_le_bind_test
\end{enumerate}
\end{proof}
